\documentclass[12pt, a4paper]{article}

% ── Packages ────────────────────────────────────────────────────────
\usepackage[utf8]{inputenc}
\usepackage[T1]{fontenc}
\usepackage[french]{babel}
\usepackage[margin=2.5cm]{geometry}
\usepackage{graphicx}
\usepackage{booktabs}
\usepackage{amsmath}
\usepackage{hyperref}
\usepackage{setspace}
\usepackage{parskip}
\usepackage{fancyhdr}
\usepackage{titlesec}
\usepackage{xcolor}
\usepackage{listings}

\definecolor{paris1blue}{RGB}{0,61,121}
\hypersetup{colorlinks=true, linkcolor=paris1blue, citecolor=paris1blue, urlcolor=paris1blue}

\lstset{
  basicstyle=\ttfamily\small,
  breaklines=true,
  frame=single,
  backgroundcolor=\color{gray!5},
  xleftmargin=1em,
  framexleftmargin=0.5em,
}

% ── Header/Footer ───────────────────────────────────────────────────
\pagestyle{fancy}
\fancyhf{}
\fancyhead[L]{\small\textcolor{paris1blue}{Analyse de Sentiment Financier}}
\fancyhead[R]{\small\textcolor{paris1blue}{Majed \& Vidal}}
\fancyfoot[C]{\thepage}
\renewcommand{\headrulewidth}{0.4pt}

\titleformat{\section}{\Large\bfseries\color{paris1blue}}{}{0em}{}
\titleformat{\subsection}{\large\bfseries\color{paris1blue!80}}{}{0em}{}
\titleformat{\subsubsection}{\normalsize\bfseries}{}{0em}{}

\onehalfspacing

% ── Metadata ────────────────────────────────────────────────────────
\title{
  \vspace{-1cm}
  {\large Universit\'e Paris 1 Panth\'eon-Sorbonne}\\
  {\large Master IASD --- Deep Learning / NLP}\\[1cm]
  {\LARGE\bfseries\color{paris1blue} Analyse de Sentiment sur Textes Financiers}\\[0.3cm]
  {\large Benchmarking Classical ML, Deep Learning et Transformers}
}
\author{
  Houssem \textsc{Majed} \and Victoria \textsc{Vidal}
}
\date{Avril 2026}

\begin{document}

\maketitle
\thispagestyle{empty}

\begin{abstract}
Ce rapport presente un projet de comparaison de six modeles de classification
de sentiment appliques a des textes financiers. Nous avons travaille sur le
FinancialPhraseBank, un jeu de donnees de 4\,846 phrases annotees par des experts
du domaine. Notre objectif etait d'evaluer trois generations d'approches NLP
(classiques, deep learning, transformers) sur un meme protocole experimental,
puis de construire un dashboard interactif de demonstration. FinBERT, un
transformer pre-entraine sur un corpus financier, obtient les meilleurs resultats
avec un macro-F1 de 87.7\%.
\end{abstract}

\newpage
\tableofcontents
\newpage

% ════════════════════════════════════════════════════════════════════
\section{Introduction}
% ════════════════════════════════════════════════════════════════════

Les marches financiers generent chaque jour un volume considerable de textes sous
forme de communiques de presse, de rapports d'analystes, de depeches Reuters ou
encore de transcriptions de conference calls. Pouvoir extraire automatiquement le
sentiment vehicule par ces textes represente un enjeu concret pour les professionnels
de la finance, que ce soit pour alimenter des strategies de trading algorithmique ou
simplement pour synthetiser l'actualite d'une entreprise.

Notre projet s'inscrit dans ce cadre. L'idee de depart etait assez simple : prendre
un jeu de donnees de reference dans le domaine, y entrainer plusieurs modeles allant
du plus basique au plus sophistique, et observer concretement comment les performances
evoluent d'une generation d'approches NLP a la suivante. Nous voulions aussi aller
au-dela du simple benchmarking en construisant un outil de demonstration complet, avec
un dashboard web permettant de tester les modeles en conditions reelles.

Le choix des trois families de modeles n'est pas arbitraire. Les approches classiques
a base de TF-IDF representent l'etat de l'art d'il y a une dizaine d'annees. Les
reseaux recurrents comme le BiLSTM ont marque la transition vers le deep learning
en NLP. Et les transformers, depuis BERT en 2018, ont fondamentalement change la
donne en introduisant le pre-entrainement massif suivi du fine-tuning sur la tache
cible.

% ════════════════════════════════════════════════════════════════════
\section{Dataset}
% ════════════════════════════════════════════════════════════════════

\subsection{FinancialPhraseBank}

Le jeu de donnees que nous avons utilise est le FinancialPhraseBank, publie par
Malo et al. en 2014 dans le Journal of the Association for Information Science and
Technology. Il contient 4\,846 phrases en anglais extraites de communiques financiers
d'entreprises europeennes, principalement finlandaises.

Chaque phrase a ete annotee par un panel de 16 experts en finance. La consigne
d'annotation etait la suivante : ``Cette phrase, du point de vue d'un investisseur,
donne-t-elle une impression positive, negative ou neutre de l'entreprise mentionnee ?''
Seules les phrases pour lesquelles une majorite d'annotateurs (plus de 50\%) s'est
accordee sur une classe sont incluses dans le dataset final.

\subsection{Distribution des classes}

La distribution des classes presente un desequilibre notable. La classe neutre domine
avec 2\,879 phrases (59.4\%), suivie de la classe positive avec 1\,363 phrases (28.1\%).
La classe negative ne represente que 604 phrases, soit 12.5\% du total. Ce desequilibre
a des consequences directes sur l'evaluation des modeles : un classificateur naif qui
predit systematiquement ``neutre'' obtiendrait deja 59\% d'accuracy, mais un macro-F1
proche de zero.

\begin{table}[h]
  \centering
  \begin{tabular}{lrr}
    \toprule
    Classe   & Nombre & Part \\
    \midrule
    Neutral  & 2\,879 & 59.4\% \\
    Positive & 1\,363 & 28.1\% \\
    Negative &    604 & 12.5\% \\
    \midrule
    Total    & 4\,846 & 100\% \\
    \bottomrule
  \end{tabular}
  \caption{Distribution des classes dans le FinancialPhraseBank.}
\end{table}

C'est precisement pour cette raison que nous avons retenu le macro-F1 comme metrique
principale. Contrairement a l'accuracy, le macro-F1 calcule la moyenne des F1-scores
de chaque classe de maniere equitable, ce qui donne autant de poids a la classe
negative (12.5\%) qu'a la classe neutre (59.4\%).

\subsection{Split des donnees}

Nous avons divise les donnees en trois ensembles avec un split stratifie (seed=42)
pour garantir la reproductibilite et une representation proportionnelle de chaque
classe dans chaque sous-ensemble.

\begin{table}[h]
  \centering
  \begin{tabular}{lrr}
    \toprule
    Split      & Samples & Part \\
    \midrule
    Train      & 3\,391  & 70\% \\
    Validation &    485  & 10\% \\
    Test       &    970  & 20\% \\
    \bottomrule
  \end{tabular}
  \caption{Repartition du split train/val/test.}
\end{table}

Ce meme split est utilise pour tous les modeles sans exception, ce qui garantit
que les comparaisons sont equitables.

% ════════════════════════════════════════════════════════════════════
\section{Modeles developpes}
% ════════════════════════════════════════════════════════════════════

\subsection{Approche classique : TF-IDF + classifieurs lineaires}

La premiere famille de modeles que nous avons implementee represente l'approche
la plus traditionnelle en NLP. Le principe est de transformer chaque texte en un
vecteur creux via TF-IDF (Term Frequency-Inverse Document Frequency), puis
d'entrainer un classifieur lineaire sur ces representations.

Nous avons utilise des unigrammes et des bigrammes, avec un vocabulaire plafonne
a 20\,000 features et une ponderation TF sublineaire. Deux classifieurs ont ete
testes : la regression logistique (avec regularisation L2 et solver lbfgs) et le
Naive Bayes multinomial (alpha=0.1).

L'avantage principal de ces modeles est leur rapidite : l'entrainement prend moins
d'une seconde sur CPU. Ils sont aussi relativement interpretables. Leur limite
fondamentale est qu'ils traitent le texte comme un sac de mots, ignorant
completement l'ordre des termes et le contexte dans lequel ils apparaissent.

\subsection{Deep Learning : BiLSTM}

Pour la deuxieme approche, nous avons implemente un reseau de neurones recurrent
bidirectionnel. L'architecture se compose d'une couche d'embedding (dimension 64,
vocabulaire de 15\,000 tokens), suivie de deux couches BiLSTM (128 unites chacune),
d'une couche dense, et enfin d'un softmax pour la classification.

L'entrainement a ete effectue sur 20 epochs avec un batch size de 32, en utilisant
l'optimiseur Adam. Nous avons mis en place un early stopping sur la loss de
validation pour eviter le sur-apprentissage. Le temps d'entrainement est d'environ
40 secondes, que ce soit sur CPU ou GPU.

Sur le papier, le BiLSTM devrait capturer l'ordre des mots et les dependances
sequentielles mieux que l'approche TF-IDF. En pratique, les resultats ont ete
decevants, comme nous le verrons dans la section resultats. Le modele n'a pas
reussi a generaliser correctement, en particulier sur la classe negative.

\subsection{Transformers : BERT, DistilBERT, FinBERT}

La troisieme et derniere famille de modeles repose sur l'architecture Transformer,
introduite par Vaswani et al. en 2017 et popularisee par BERT (Devlin et al., 2019).
Le principe est de partir d'un modele pre-entraine sur un tres large corpus de
textes generiques, puis de le fine-tuner sur notre tache specifique.

Nous avons fine-tune trois variantes. BERT (``base-uncased'') est le modele original
avec 110 millions de parametres, pre-entraine sur BookCorpus et Wikipedia. DistilBERT
est une version distillee de BERT, 40\% plus legere (66 millions de parametres) tout
en conservant environ 97\% de ses performances sur les benchmarks generiques. FinBERT
est une version de BERT qui a subi un pre-entrainement supplementaire sur un corpus
de textes financiers (rapports d'analystes, earnings calls, communiques de presse).
Ce pre-entrainement additionnel lui permet de comprendre les nuances propres au
vocabulaire financier, comme les termes ``restructuring'', ``write-down'' ou
``impairment charge'' qui ont une connotation negative specifique au domaine.

Pour les trois modeles, les hyperparametres de fine-tuning sont identiques : 5 epochs,
learning rate de 2e-5, avec early stopping sur la performance de validation. Nous
avons du mettre en place plusieurs optimisations GPU pour faire tourner l'entrainement
sur notre carte graphique (RTX 5070 Laptop, 8 Go de VRAM) : precision mixte fp16,
gradient checkpointing pour BERT et FinBERT (qui sont plus gourmands en memoire),
gradient accumulation, et un plafonnement de l'utilisation VRAM a 85\% pour eviter
les erreurs de memoire.

% ════════════════════════════════════════════════════════════════════
\section{Resultats}
% ════════════════════════════════════════════════════════════════════

\subsection{Tableau comparatif}

\begin{table}[h]
  \centering
  \begin{tabular}{clcccc}
    \toprule
    Rang & Modele & Accuracy & Macro-F1 & Type & Temps \\
    \midrule
    1 & \textbf{FinBERT}    & 88.1\% & \textbf{87.7\%} & Transformer   & 82s  \\
    2 & DistilBERT           & 86.8\% & 86.0\%           & Transformer   & 51s  \\
    3 & BERT                 & 85.5\% & 83.8\%           & Transformer   & 111s \\
    4 & TF-IDF + NaiveBayes  & 71.3\% & 62.7\%           & Classical     & <1s  \\
    5 & TF-IDF + LogReg      & 72.4\% & 60.1\%           & Classical     & <1s  \\
    6 & BiLSTM               & 63.7\% & 37.3\%           & Deep Learning & 41s  \\
    \bottomrule
  \end{tabular}
  \caption{Benchmark des 6 modeles sur le split de test (970 exemples).}
\end{table}

\subsection{Analyse des resultats}

Les resultats montrent une hierarchie assez nette. FinBERT arrive en tete avec un
macro-F1 de 87.7\%, ce qui confirme l'apport du pre-entrainement specifique au
domaine financier. L'ecart avec BERT generique (+3.9 points de macro-F1) est
significatif et illustre bien l'interet d'adapter le modele au vocabulaire cible.

Un resultat surprenant est la performance de DistilBERT (86.0\%), qui depasse BERT
complet (83.8\%) malgre un nombre de parametres inferieur de 40\%. Notre hypothese
est que sur un petit dataset comme le FinancialPhraseBank, un modele plus compact
est moins susceptible d'overfitter. DistilBERT est aussi nettement plus rapide a
fine-tuner (51s contre 111s pour BERT).

Les baselines classiques montrent un phenomene interessant : leur accuracy est
correcte (71-72\%) mais leur macro-F1 est bien plus faible (60-63\%). Cet ecart
s'explique par le fait que ces modeles predisent majoritairement ``neutre'', ce qui
leur donne une bonne accuracy grace a la prevalence de cette classe, mais les fait
echouer sur les classes minoritaires.

Le cas le plus parlant est celui du BiLSTM, qui obtient le pire macro-F1 du benchmark
(37.3\%). En analysant les predictions, nous avons constate qu'il obtenait un recall
de 0\% sur la classe negative. Sans pre-entrainement et avec seulement 5\,000
exemples d'entrainement, le reseau recurrent n'a tout simplement pas assez de
donnees pour apprendre les patterns de chaque classe, surtout la classe minoritaire.

\subsection{Enseignements}

Le principal enseignement de ce benchmark est le role crucial du pre-entrainement
dans un contexte de faible volume de donnees. Les transformers pre-entraines
surpassent largement les approches non pre-entrainees, et le pre-entrainement
specialise (FinBERT) apporte un gain supplementaire. Ce constat est coherent avec
la litterature recente en NLP : le paradigme ``pre-train then fine-tune'' est
particulierement efficace quand le dataset cible est petit.

% ════════════════════════════════════════════════════════════════════
\section{Dashboard et fonctionnalites}
% ════════════════════════════════════════════════════════════════════

\subsection{Architecture du dashboard}

Pour rendre nos modeles accessibles et demonstrables, nous avons developpe un
dashboard web en utilisant Streamlit. L'interface se compose de plusieurs pages
accessibles via un menu de navigation lateral. Le design a ete entierement
personnalise avec un systeme de CSS injecte, comprenant une sidebar sombre, des
cartes avec bordures et ombres, et une typographie Inter pour un rendu professionnel.

Le modele actif est selectionnable directement depuis la page d'analyse, a cote
des boutons d'action. Le chargement des modeles est mis en cache grace a
\texttt{st.cache\_resource}, ce qui evite de recharger les poids a chaque interaction.

\subsection{Analyse individuelle et comparaison multi-modeles}

La page principale permet deux modes d'utilisation. Le premier est l'analyse
individuelle : l'utilisateur saisit ou dicte une phrase, choisit un modele, et
obtient la prediction avec les scores de confiance pour chaque classe. Le second
mode est la comparaison multi-modeles, qui fait tourner la meme phrase dans les
trois transformers disponibles et affiche les resultats cote a cote. Ce mode
inclut egalement une analyse LIME (Local Interpretable Model-Agnostic Explanations)
pour chaque modele, ce qui permet de voir quels mots ont le plus influence la
prediction de chacun.

\subsection{Batch Analysis}

La page Batch Analysis permet de traiter un fichier CSV complet. L'utilisateur
uploade un fichier contenant une colonne de headlines, et le systeme classifie
chaque ligne avec le modele actif. Les resultats sont presentes sous forme de
metriques synthetiques (nombre par classe), de graphiques (camembert de distribution,
histogramme de confiance), d'un tableau detaille, et d'un bouton de telechargement
CSV.

\subsection{Live News via Yahoo Finance}

Pour aller au-dela des donnees statiques, nous avons integre une fonctionnalite
d'analyse en temps reel. L'utilisateur entre un ticker boursier (par exemple AAPL
pour Apple) et le systeme utilise la librairie \texttt{yfinance} pour recuperer les
dernieres headlines d'actualite associees a ce titre. Chaque headline est ensuite
classifiee par le modele actif, et un score de sentiment global est calcule en
faisant la difference entre la proportion de headlines positives et negatives.

Le principal defi technique a ete l'extraction des titres depuis la structure de
donnees renvoyee par l'API yfinance, qui varie selon les versions et les tickers.
Nous avons implemente un parsing robuste qui gere les differents formats possibles.

\subsection{Entree vocale}

La derniere fonctionnalite notable est l'entree vocale. L'idee etait de permettre
a l'utilisateur de dicter une headline plutot que de la taper, a la maniere des
interfaces conversationnelles modernes. L'implementation repose sur la Web Speech
API du navigateur (compatible Chrome et Edge).

Le defi principal etait lie a l'architecture de Streamlit : tous les composants
HTML personnalises sont rendus dans des iframes sandboxees qui n'ont pas acces au
microphone du navigateur. Notre premiere tentative avec \texttt{streamlit-javascript}
a echoue car le code s'executait dans le contexte de l'iframe. La solution finale
utilise \texttt{components.html()} comme un simple bootstrap pour injecter un
element \texttt{<script>} dans le document parent de Streamlit. Ce script s'execute
dans le contexte de la page principale, ce qui lui donne acces au microphone et
permet de satisfaire l'exigence de ``user gesture'' de Chrome. La transcription est
ensuite injectee dans le textarea via le setter React natif de l'element HTML, suivi
d'evenements synthetiques pour que Streamlit prenne en compte la nouvelle valeur.

% ════════════════════════════════════════════════════════════════════
\section{Difficultes rencontrees}
% ════════════════════════════════════════════════════════════════════

Tout au long du projet, nous avons fait face a un certain nombre de difficultes
techniques que nous souhaitons documenter ici.

La premiere difficulte concernait la gestion de la memoire GPU. Le fine-tuning de
BERT et FinBERT (110 millions de parametres chacun) sur une carte avec 8 Go de
VRAM necessitait des optimisations specifiques. Nous avons du activer la precision
mixte fp16, le gradient checkpointing (qui echange du temps de calcul contre de la
memoire), la gradient accumulation pour simuler des batch sizes plus grands, et un
plafonnement de l'allocation memoire a 85\% de la VRAM disponible. Sans ces mesures,
l'entrainement echouait systematiquement avec des erreurs CUDA out-of-memory.

La deuxieme difficulte etait liee a l'installation de PyTorch avec support CUDA
sur Windows. La commande standard \texttt{pip install torch} installe une version
CPU-only. Il a fallu configurer un index de paquets specifique (\texttt{pytorch-cu128})
dans le fichier \texttt{pyproject.toml} via la directive \texttt{[tool.uv.sources]}
pour obtenir une version compatible avec notre driver CUDA.

La troisieme difficulte, et probablement la plus chronophage, a ete l'implementation
de l'entree vocale dans Streamlit. Nous avons explore et abandonne plusieurs pistes
avant de trouver la solution d'injection dans le contexte parent. Le probleme est
que les iframes Streamlit ont un attribut \texttt{sandbox} qui bloque l'acces au
microphone, et que les permissions du navigateur exigent que la creation du
\texttt{SpeechRecognition} se fasse dans le bon contexte JavaScript avec un geste
utilisateur explicite.

Enfin, la synchronisation entre les modifications JavaScript du DOM et l'etat interne
de Streamlit a pose des problemes subtils. Quand on modifie la valeur d'un textarea
via JavaScript, Streamlit ne detecte pas automatiquement le changement. Nous avons du
utiliser le setter natif du prototype \texttt{HTMLTextAreaElement}, dispatcher des
evenements synthetiques \texttt{input} et \texttt{change}, et utiliser des callbacks
\texttt{on\_click} plutot que des affectations directes au \texttt{session\_state}
pour contourner les restrictions de Streamlit sur la modification de l'etat apres
l'instanciation d'un widget.

% ════════════════════════════════════════════════════════════════════
\section{Conclusion}
% ════════════════════════════════════════════════════════════════════

Ce projet nous a permis de concretiser et comparer trois generations d'approches
NLP sur une tache reelle de classification de sentiment financier. Le resultat
principal est sans ambiguite : FinBERT domine le benchmark avec un macro-F1 de
87.7\%, confirmant que le pre-entrainement specifique au domaine est le facteur
le plus determinant pour la performance sur un petit dataset specialise.

Au-dela du benchmarking, le developpement du dashboard Streamlit a ete l'occasion
de confronter nos modeles a des cas d'usage concrets (analyse en temps reel,
traitement par lot, explicabilite) et de resoudre des problemes d'ingenierie
logicielle non triviaux (gestion GPU, integration Web Speech API, synchronisation
JavaScript-Python).

Parmi les pistes d'amelioration, on peut envisager l'elargissement du dataset
d'entrainement (SEC filings, earnings call transcripts), l'experimentation avec
des modeles plus recents comme DeBERTa ou des approches few-shot a base de LLMs,
et le deploiement du dashboard sur une infrastructure cloud avec une API REST
pour un acces plus large.

% ════════════════════════════════════════════════════════════════════
\section*{References}
% ════════════════════════════════════════════════════════════════════

\begin{enumerate}
  \item Malo, P., Sinha, A., Korhonen, P., Wallenius, J., \& Takala, P. (2014).
    Good debt or bad debt: Detecting semantic orientations in economic texts.
    \emph{Journal of the Association for Information Science and Technology},
    65(4), 782--796.
  \item Araci, D. (2019). FinBERT: Financial Sentiment Analysis with Pre-trained
    Language Models. \emph{arXiv preprint arXiv:1908.10063}.
  \item Devlin, J., Chang, M.-W., Lee, K., \& Toutanova, K. (2019). BERT:
    Pre-training of Deep Bidirectional Transformers for Language Understanding.
    \emph{Proceedings of NAACL-HLT}.
  \item Vaswani, A., Shazeer, N., Parmar, N., et al. (2017). Attention Is All
    You Need. \emph{Advances in Neural Information Processing Systems (NeurIPS)}.
  \item Sanh, V., Debut, L., Chaumond, J., \& Wolf, T. (2019). DistilBERT, a
    distilled version of BERT: smaller, faster, cheaper and lighter.
    \emph{arXiv preprint arXiv:1910.01108}.
  \item Ribeiro, M. T., Singh, S., \& Guestrin, C. (2016). ``Why Should I Trust
    You?'': Explaining the Predictions of Any Classifier. \emph{KDD}.
\end{enumerate}

\end{document}
